% TeX root=../main.tex

\section{Situação linguística em sincronia}

\subsection{Posição do eslavo antigo no conjunto eslávico}
O eslavo antigo pertence ao grupo das línguas eslavas
meridionais (\emph{cf.}~\autoref{fig:reino_bulgaria}).
As línguas eslavas são divididas em três grandes
grupos, o eslavo meridional (búlgaro, macedônio, servo-croata e esloveno), o
eslavo oriental (russo, ucraniano e bielorrusso) e o eslavo ocidental (\glsxtrlong{checo},
eslovaco, sorábio, polonês, cachubo e eslovíncio, os dois últimos já
extintos). Esses três grandes grupos advêm de uma divisão do eslavo comum, um
ramo do indo-europeu. Entre o indo-europeu e o eslavo comum seria verossímil
reconstruir uma unidade intermediária chamada de balto-eslavo, que seria a
origem do grupo eslavo de um lado e do grupo báltico de outro, mas essa tese
não é aceita por todos. A divisão do eslavo comum é consequência da expansão
geográfica do eslavo nos séculos 5 e 6 EC: a partir de um território inicial
que corresponde largamente ao norte da Ucrânia atual e ao sul da Polônia, os
Eslavos avançaram em direção ao oeste em território germânico, ao sul, onde se
instalaram nos Bálcãs, e alguns chegaram até mesmo ao Peloponeso; e,
paralelamente, colonizaram o noroeste da Rússia. O eslavo antigo, cujos primeiros
registros escritos remontam ao séc. 9, está, por isso, ainda próximo do eslavo
comum. Ainda assim, ele pertence já distintamente ao grupo meridional e, mais
especificamente, ao subgrupo oriental do eslavo meridional, representado pelo
búlgaro e pelo macedônio. Com efeito, há traços fonéticos presentes no eslavo antigo
característicos desse subgrupo que o opõem ao subgrupo ocidental, representado
pelo sérvio-croata e esloveno. Mesmo assim, essas divisões do grupo eslavo
ainda não eram tais que impedissem a intercompreensão. De fato, Cirilo e
Método falavam um dialeto eslavo meridional e traduziram o Evangelho para seu
dialeto, mas o destinatário da tradução era a população eslava da Morávia, que
falava um dialeto eslavo ocidental. Não havia necessidade de uma grande
adaptação para garantir a compreensão.

O eslavo antigo compartilha as características das línguas indo-europeias: é uma língua
flexional (a flexão nominal conta sete casos) do tipo sintético, com uma rica
morfologia derivacional usando afixos variados com ordem de palavras livre. 

\subsection{Contatos linguísticos}
O subgrupo oriental do eslavo meridional (búlgaro-macedônio) se instalou em um
território ocupado já há muito tempo pelos trácios. Os trácios foram
eslavizados e a questão de um substrato trácio não está posta em relação ao
eslavo antigo, mas a ausência de conhecimento sobre o trácio faz com que nos seja
impossível ver a ação de um tal substrato. Posteriormente, a chegada dos
búlgaros, que impõem sua dominação e fundam o primeiro reino da Bulgária
(681), introduz nessa região falantes de uma língua não indo-europeia, mas os
búlgaros também foram rapidamente eslavizados e não é possível falar de uma
real influência de superstrato proto-búlgaro sobre o eslavo antigo.

O eslavo antigo confinava também com outras línguas indo-europeias, o romance da Panônia
(que se tornará o romeno) ao norte e o grego ao sul. Se a influência do
romance é quase indetectável, a influência do grego é manifesta, tanto no
léxico como na sintaxe. De qualquer forma, a influência grega é essencialmente
um fato da tradução, dado que o eslavo antigo calca o original grego, e é muito difícil
saber em que medida a própria língua foi influenciada pelo grego. É evidente
que, visto que o grego era a língua de cultura e o eslavo a língua vernacular,
o estatuto respectivo das duas línguas fazia que o uso linguístico das elites
letradas, helenizadas, devia certamente apresentar traços remontáveis ao
grego.

\subsection{Dialetos}
O eslavo nos séculos 9 e 10 tinha diferenças dialetais. É possível distinguir
nos textos de eslavo antigo que nos chegaram dois ramos, um do norte, menos bem
atestado, o do eslavão morávio, e um ramo do sul, o melhor atestado, que se
divide entre um dialeto búlgaro e um macedônio. A língua de Cirilo e Método
era seu dialeto natal, o de Tessalônica, ao sul do território e, portanto, um
dialeto macedônio. Contudo, os textos em eslavo antigo não conservaram a homogeneidade
linguística inicial: os textos traduzidos no contexto da missão morávia
(863-885) por discípulos formados pelos missionários foram influenciados pelo
dialeto local, que é um dialeto eslavo ocidental, não meridional; eles
comportam, portanto, um certo número de traços eslavos ocidentais (fala-se de
moravismos) sobre o pano de fundo do eslavo meridional, como se testemunha no
\glslink{Kiev}{\emph{Missal}}, que é nosso único testemunho da tradição morávia.
Consequentemente, o retorno à Bulgária dos discípulos de Método expulsos da
Morávia marcou o início de um processo de ``bulgarização'' da língua. Os
manuscritos que chegaram a nós datam todos ou quase todos do período búlgaro
ou macedônio, entre o fim do séc. 10 e o fim do 11, e refletem certa
bulgarização. Eles se repartem em dois grandes grupos segundo representam a
tradição búlgara da escola de Preslav ou a tradição macedônica da escola de
Ócrida, nova capital de Samuel. No conjunto, a escola de Preslav é mais
influenciada pelo grego em sua prática tradutória que a de Ócrida, que é mais
conservadora e sem dúvida mais fiel à tradição cirilo\hyp{}metodiana. Isso se
reflete também no alfabeto escolhido: a escola de Preslav adotou o alfabeto
cirílico enquanto a escola de Ócrida manteve o glagolítico. 

As principais divergências entre o dialeto morávio (eslavo ocidental) e os 
dialetos búlgaro e macedônio (eslavo meridional), são as seguintes 
(notar-se-á que todos esses pontos não estão atestados no \glslink{Kiev}{\emph{Missal}}, do
qual nos restam apenas sete fólios reto-verso):

\begin{compactitem}[--]
  \item resultado \emph{š} a partir de \emph{x} na segunda e terceira
      palatalizações, em vez do \emph{s} do eslavo meridional (sem exemplo no
      \glsxtrlong{Kiev}); manutenção de \ocsex{kvE}, \ocsex{gvE} em comparação
      com \ocsex{cvE}, \ocsex{dzvE} em eslavo meridional (\glsxtrlong{checo}
      \ocsex{cvEt} ``flor'' em comparação como o eslavo antigo \ocsex{cvEtъ},
      sem exemplo no \glslink{Kiev}{\emph{Missal}});
	\item manutenção dos grupos \emph{tl}, \emph{dl}, simplificados em \emph{l} no
    eslavo meridional (particípio perfeito \glsxtrlong{checo} \emph{vedl}
    ``tendo conduzido'' em comparação com o eslavo antigo \emph{velъ}, sem
    exemplo no \glslink{Kiev}{\emph{Missal}});
	\item resultado de *\emph{\#RăC} > \emph{\#RoC} a partir de *\emph{\#ăRC}
    original (metátese sem alongamento compensatório) em comparação com
    *\emph{\#RāC} > \emph{\#RāC} em eslavo meridional (metátese com alongamento
    compensatório), não atestado no \glslink{Kiev}{\emph{Missal}}, mas é
    possível encontrar nos manuscritos meridionais a forma
    \cirilocsex{робъ,robъ} ``servidor/escravo'' em vez de
    \cirilocsex{рабъ,rabъ}, cujo pervérbio \cirilocsex{роз-,roz-}, indicando a
    separação, em vez de \cirilocsex{раз-,raz-}; as formas em \ocsex{ro-} são
    moravismos; inversamente, \ocsex{razdrEšen’ie} ``remissão'',
    \passageme{Kiev5}, é uma forma de origem meridional;
	\item resultado \emph{c}, \emph{z}, *\emph{tj}, *\emph{dj}, em vez de
    \ocsex{St}, \ocsex{Zd} em búlgaro-macedônio:
    \ocsex{nasyceni}, \passageme{Kiev6223}, em vez de \ocsex{nasySteni} ``saciar-se'', e a título
    de moravismo em exemplos isolados como o \gloss{gen.sg.neut.} \ocsex{rozъstva}
    (\passageme{Mt146}
    \glsxtrshort{Marianus}) por \ocsex{roZdIstva} ``nascimento''.
\end{compactitem}

Dentro do ramo meridional, as principais divergências entre o dialeto búlgaro antigo e o dialeto macedônio antigo são as seguintes.

Foneticamente, o \emph{jer} anterior /\jeru/ é uma vogal arredondada em
macedônio antigo, mas não arredondada no búlgaro antigo, como se vê pelo
tratamento do /\jeru/ forte, às vezes grafado \grafema{о} nos textos macedônios
antigos, mas sempre \grafema{ъ} nos textos búlgaros antigos.
O resultado \ocsex{ja} > \ocsex{E} caracteriza os textos de proveniência
macedônia, enquanto ja é conservado nos textos de origem búlgara; essa última
diferença é antiga, uma vez que o alfabeto glagolítico tem apenas um grafema
para /ja/ e /\ocstrans{E}/, o que demonstra que a confusão já tinha sido adquirida na
segunda metade do séc. 9 no momento da criação do alfabeto glagolítico.
O /r’/ tem a tendência de se endurecer em búlgaro (\cirilocsex{цѣсара,cEsara}
por \cirilocsex{цѣсара,cEsar’a} ``César'' \glsxtrshort{Supr}), enquanto em
macedônio antigo ele permanece palatal.
O /dz/ africado é bem conservado em macedônio antigo mas já é
simplificado em /z/ em búlgaro antigo, como é o caso do eslavo ocidental
(\glsxtrshort{Kiev}).
O [l] epentético resultado da palatalização das labiais
tende a ser eliminado em fronteira morfêmica nos manuscritos búlgaros antigos
(\ocsex{na zemi} ``sobre a terra'' \glsxtrshort{Supr}, em vez de \ocsex{na
zeml’i}), mas esse desenvolvimento é encontrado também no eslavo ocidental, e se
os manuscritos macedônios antigos não têm traço algum disso, ou quase nenhum, é
sobretudo por seu caráter conservador que lhes faz respeitar escrupulosamente a
grafia antiga. 

Quanto à morfologia nominal, o búlgaro antigo conserva no conjunto das
desinências dos temas em \emph{-i-} (\gloss{inst.sg masc.} -\ocsex{ьmь},
\gloss{dat.pl.} -\ocsex{ьmъ}, \gloss{loc.pl.} -\ocsex{ьхъ}), enquanto o
macedônio antigo tende a lhes substituir pelos alomorfes -\ocsex{emь},
-\ocsex{emъ}, -\ocsex{exъ} derivados dos temas em \ocsex{-jo-}.

Quanto à morfologia verbal, a forma de \gloss{3ªdual} \ocsex{-te} é conservada
em macedônio antigo, mas é substituída em búlgaro antigo por \ocsex{-ta},
originalmente a desinência de \gloss{2ªdual}.
O aoristo sigmático de tipo \ocsex{vEsъ} é bem conservado em macedônio antigo,
mas é quase completamente eliminado em búlgaro antigo em face da nova forma
\ocsex{vedoxъ}.
Quanto ao léxico, algumas divergências podem ser notadas.
Há no dialeto búlgaro algumas formas de origem turca (empréstimos do
proto-búlgaro) que estão ausentes no dialeto macedônio.
Além dos empréstimos, essas divergências contrapõem tanto dois lexemas
diferentes quanto duas derivações diferentes de uma mesma raiz
(\autoref{tab:diverglex}).


\begin{table}[!htb]
  \caption{Divergências lexicais entre os dialetos eslavos
  antigos}\label{tab:diverglex}
  \begin{center}
    \begin{tabular}[c]{lll}
     \toprule 
      Significado & Macedônio antigo &	Búlgaro antigo\\
      \midrule
``combate'' & \cirilocsex{брань,branь}  & \cirilocsex{рать,ratь}\\
``tempo, hora'' & \cirilocsex{година,godina}  & \cirilocsex{часъ,časъ}\\
``salário'' & \cirilocsex{мꙑто,myto}  & \cirilocsex{мьзда,mьzda}\\
``ilha'' & \cirilocsex{отокъ,otokъ}  & \cirilocsex{островъ,ostrovъ}\\
``pastor'' & \cirilocsex{pastyr'I}  & \cirilocsex{пастѹхъ,pastuxъ}\\
``alimentar'' & \cirilocsex{питати,pitati} & \cirilocsex{кръмити,krъmiti}\\
``por causa de'' & \cirilocsex{ради,radi}  & \cirilocsex{dEl'ja,dEl’a}\\
``reunida'' & \cirilocsex{съньмъ,sъnьmъ}  & \cirilocsex{съборъ,sъborъ}\\
``somente'' & \cirilocsex{тъкъмо,tъkъmo}  & \cirilocsex{тъчиѭ,tъčijq}\\
``esquerda''	 & \cirilocsex{шѹи,šujь}  & \cirilocsex{лѣвъ,lEvъ}\\
\bottomrule
    \end{tabular}
  \end{center}
\end{table}


No que concerne à situação dialetal, é possível ver diretamente dois dialetos no
eslavo meridional (oriental, ``búlgaro'', e ocidental, ``macedônio'') e um
indiretamente: \glsxtrshort{Marianus} apresenta de fato traços de tipo sérvio,
notadamente a desnasalização de \ocsex{õ} em \ocsex{u} (que é regular em sérvio
mas não é atestada em macedônio e búlgaro), o que faz supor que o monge que
copiou o manuscrito era falante de um dialeto de transição entre macedônio e
sérvio (\emph{cf.}~\ref{sec:invvoc}).

\subsection{Traços balcânicos}
O búlgaro e o macedônio modernos pertencem ao grupo das línguas balcânicas, que
são caracterizadas por uma convergência areal bastante forte.
Os principais traços das línguas balcânicas são a perda do infinitivo,
substituído por subjuntivos, o sincretismo do genitivo e do dativo (que em
búlgaro-macedônio moderno resulta na perda quase completa da flexão nominal), a
formação do futuro a partir de uma perífrase volitiva (``eu quero fazer'' por
``eu farei''), a existência de um perfeito perifrástico com auxiliar ``haver''
e, em algumas línguas, o desenvolvimento de um artigo posposto.
Nenhum desses traços pode ser claramente observado em eslavo antigo.
De modo geral, é possível detectar alguns traços esporádicos, o que se explica
sem dúvida sobretudo pelo fato de que o eslavo antigo é uma língua de tradução
para uso religioso, de que o original grego não apresenta esses traços
balcânicos, dado que é bem mais antigo, e de que se esses desenvolvimentos
existiam na língua popular, eles eram provavelmente característicos de um nível
da língua evitada pela tradução de textos sagrados. 

\section{História da língua}

\subsection{Consonantismo}\label{sec:consoantes}
O ramo eslavo, como também o báltico, combinou as sonoras aspiradas originais
com sonoras não aspiradas, de modo que não há mais que dois modos de articulação
para as consoantes, surda e sonora: \cirilocsex{dati} ``dar'' < *\emph{dō}- <
*\pietrans{deh3-} (\glsxtrshort{grego}~δίδωμι, \glsxtrshort{latim}~\emph{do},
\glsxtrshort{sânscrito}~\emph{dā-}),  \cirilocsex{dElo} ``ato'' < *\emph{dhē}- <
*\pietrans{dheh1-} (\glsxtrshort{grego}~τίθημι,
\glsxtrshort{latim}~\emph{facio}, \glsxtrshort{inglês}~\emph{do},
\glsxtrshort{sânscrito}~\emph{dhā-}).
De qualquer forma, há um traço da diferença etimológica entre sonora e sonora
aspirada pela qual uma vogal breve que precede uma antiga sonora não aspirada é
refletida em eslavo e em báltico por uma vogal longa de entonação rude (lei de
Winter, que, contudo, não é aceita por todos), enquanto uma vogal breve que
precede uma antiga sonora não aspirada continua breve.
É a lei de Winter que explica as formas com *\pietrans{E} > \ocsex{E} nas raízes
*\pietrans{sed-} ``estar sentado'', *\pietrans{h1ed-} ``comer'', entre outras.

\begin{table}[!htb]
  \caption{Sistema consonantal do indo-europeu ao eslavo antigo}\label{tab:pieeslavcons}
  \begin{center}
    \begin{tabular}[c]{lll}
      \toprule
*\emph{p} > \emph{p} & &*\pietrans{bh} > \emph{b}\\
\midrule
*\emph{t} > \emph{t} &*\emph{d} > \emph{d} &*\pietrans{dh} > \ocsex{d}\\
\midrule
*\pietrans{c} > \emph{s} & *\pietrans{j} > \ocsex{z} &*\pietrans{jh} > \ocsex{z}\\
\midrule
*\emph{k} > \emph{k/č} &*\emph{g} > \emph{g/ž} &*\pietrans{gh} > \ocsex{g/ž}\\
\midrule
*\pietrans{kw} > \emph{k/č} &*\pietrans{gw} > \emph{g/ž} &*\pietrans{gwh} > \emph{g/ž}\\
\midrule
\multicolumn{3}{l}{*\emph{r}, *\emph{l}, *\emph{n}, *\emph{m} > \emph{r},
\emph{l}, \emph{n}, \emph{m}}\\
\midrule
\multicolumn{3}{l}{*\emph{s} > \emph{s/š/x} }\\
\multicolumn{3}{l}{de acordo com o contexto (RUKI)}\\
\midrule
\multicolumn{3}{l}{*\pietrans{y}̯, *\pietrans{w} > \ocsex{j}, \ocsex{v}}\\
\bottomrule
\multicolumn{3}{l}{}\\
\multicolumn{3}{l}{NB\@: o indo-europeu não possuía /b/.}\\
    \end{tabular}
  \end{center}
\end{table}


O eslavo pertence ao grupo de línguas indo-europeias \emph{satem}: as palatais
originais resultam em sibilantes (*\pietrans{c} > \emph{s}:
\cirilocsex{срьдьце,srьdьce} ``coração'', de *\pietrans{cRdi-ko-},
\glsxtrshort{latim}~\emph{cor, cordis}, \glsxtrshort{grego}~καρδία,
\glsxtrshort{inglês}~\emph{heart}; *\pietrans{j} e *\pietrans{jh} > \emph{z}:
\cirilocsex{зима,zima} ``inverno'', de *\pietrans{jheim-},
\glsxtrshort{latim}~\emph{hiems}, \glsxtrshort{grego}~χειμών) e as labiovelares
originais e velares próprias resultam em velares (*\pietrans{kw} > \emph{k}:
interrogativo \cirilocsex{которꙑи,kotoryjь} ``qual?'', de *\pietrans{kwo-tero-},
\glsxtrshort{grego}~πότερος, \glsxtrshort{latim}~\emph{uter},
\glsxtrshort{sânscrito}~\emph{katara-}; *\pietrans{gw} *\pietrans{gwh} >
\emph{g}: \cirilocsex{гонити,goniti} ``caçar'', de *\pietrans{gwhen-},
\glsxtrshort{grego}~φόνος, \glsxtrshort{sânscrito}~\emph{ghana-}).
A tendência à palatalização é mantida: após a palatalização \emph{satem}, o
eslavo ainda teve três outras palatalizações.
A primeira, que afeta as velares (i.e., originalmente labiovelares ou velares
próprias) diante de *\emph{i}, *\emph{e}, resultou em africadas chiadas: o
caráter africado é conservado na surda e eliminado na sonora (*\pietrans{kwi},
*\pietrans{kwe} > *\emph{ki}, *\emph{ke} > \emph{či, če}:
\cirilocsex{чєтꙑрє,četyre} ``quatro'', de *\pietrans{kwetūres},
\glsxtrshort{grego}~τέτταρες, \glsxtrshort{latim}~\emph{quattuor},
\glsxtrshort{sânscrito}~\emph{catvāraḥ}; *\pietrans{gwi}, *\pietrans{gwe}
*\pietrans{gwhi}, *\pietrans{gwhe} > *gi, *ge > *dži, *dže > ži, že:
\cirilex{жєна} ``mulher'', de *\pietrans{gwenA}, \glsxtrshort{grego}~γυνή,
\glsxtrshort{inglês}~queen ``rainha'').
A segunda, que afeta também velares, mas as resultadas de um E derivado da
monotongação de *\emph{ai} (<*\emph{oi}, *\emph{ai}), produz africadas
sibilantes (*\ocsex{kE} > \ocsex{cE}: \cirilocsex{цѣна,cEna} ``preço''
*\pietrans{kwoynA}, \glsxtrshort{lituano}~\emph{kainà},
\glsxtrshort{grego}~ποινή ``preço, pena''; *\ocsex{gE} > \ocsex{dzE}, o último
já desafricado em \ocsex{zE} em eslavo antigo: encontra-se tanto a grafia
\cirilocsex{ꙁѣло,dzElo} como \cirilocsex{зѣло,zElo} ``muito'' < *\emph{goilo-},
\glsxtrshort{lituano}~\emph{gailùs} ``brusco, violento'').
A terceira, que é uma palatalização progressiva e não regressiva como as
anteriores (a velar é palatalizada por um *\emph{i} precedente:
\cirilocsex{срьдьцє,srьdьce} ``coração'' < *\pietrans{cRdi-ko-}), tem os mesmos
resultados que a segunda, ou seja, \emph{c} para a surda e \emph{dz} pela
sonora.
Enfim, toda consoante diante de *\emph{j} se palataliza igualmente.
O eslavo é, portanto, dotado de um grande número de fonemas palatais
(\emph{cf.}~\autoref{tab:palatalizações}).

Essas palatalizações são a causa direta de numerosas alterações morfológicas
(\emph{cf.}~\ref{sec:altmorfofon}).
Em línguas eslavas modernas, as alterações são bem conservadas na
flexão verbal e na derivação, menos bem na flexão nominal. 

\begin{table}[!htb]
  \caption{Palatalizações do eslavo comum para o eslavo
  antigo}\label{tab:palatalizações}
  \begin{center}
    \begin{tabular}[c]{llll}
      \toprule
       & *tj & \ocsex{št} & \cirilex{шт}\\
      Dentais & *dj & \ocsex{žd} & \cirilex{жд}\\
              & *sj & \ocsex{s}̌ & \cirilex{ш}\\
      \midrule
      Labiais & *pj, *bj, *mj & \ocsex{pl'}, \ocsex{bl’}, \ocsex{ml’} & \cirilex{пл}, \cirilex{бл}, \cirilex{мл}\\
      \midrule
      Soantes & *nj, *rj, *lj & \ocsex{n'}, \ocsex{r’}, \ocsex{l’} & \cirilex{н},\cirilex{р}, \cirilex{л}\\
      \midrule
      Velares: & *kj, *kī̆, *kē̆ & \ocsex{C} & \cirilex{ч}\\
      primeira palatalização & *gj, *gī̆, *gē̆ & \ocsex{Z} & \cirilex{ж}\\
      \midrule
      Velares: & *kě (*kai) &	\ocsex{c} & \cirilex{ц}\\
      segunda palatalização	& *\ocstrans{gE} & \ocsex{dz/z} &	\cirilex{ꙁ}\slash{}\cirilex{з}\\
                            & *\ocsex{xE} & \ocsex{s} (\ocsex{š} em morávio) & \cirilex{с}\\
      \bottomrule
    \end{tabular}
  \end{center}
\end{table}

Como o indo-iraniano e outras línguas do grupo \emph{satəm}, o eslavo tem
palatalização de *\emph{s} após *\emph{r}, *\emph{k}, *\emph{u}, *\emph{i}
(regra RUKI).
O resultado é duplo: *\emph{s} > \emph{š} como em indo-iraniano diante de vogal
anterior (\gloss{\gloss{2sg.}} dos verbos em \emph{-i-} eslavo antigo \emph{-i-ši} <
*\emph{ei-sēi}; a variante \emph{-ši} foi estendida para todos os tipos verbais,
por isso há em eslavo antigo \emph{-e-ši} em vez do *\emph{-e-si} esperado);
*\emph{s} > \emph{x} diante de vogal posterior (\gloss{loc.pl.} dos temas em
-\emph{o}-, eslavo antigo \ocsex{-Exъ} < \mbox{*-\emph{oixŭ}} < *-\emph{oisu},
\glsxtrshort{sânscrito}~-\emph{eṣu}, \glsxtrshort{grego}~homérico~-\emph{oisi};
a variante \ocsex{x} é ainda estendida a todos os tipos nominais, por isso os temas em
-\emph{a}- do eslavo antigo apresentam -\ocsex{a-xъ} em vez da forma esperada
*-\ocsex{a-sъ}); diante de consoante, por fim, *\emph{s} não muda.
Ainda assim, essa evolução teve por consequência a aparição de variantes para
certos morfemas (assim o aoristo \cirilocsex{бꙑти,byti} ``ser'' <*\pietrans{bhuH}-: \gloss{\gloss{1sg.}.}
\cirilocsex{бꙑхъ,byxъ}, \gloss{3sg.} \cirilocsex{бꙑстъ,bystъ}, \gloss{3pl.}
\cirilocsex{бꙑшѧ,byšẽ}: /x/, /s/ e /š/ correspondem ao morfema *-\emph{s}- de
  aoristo em três diferentes contextos).

As líquidas e nasais são conservadas sem alteração, exceto em posição final de
sílaba.
Um *\emph{m} se torna \emph{n} em final de palavra, dado que o eslavo era uma
língua de n final como o báltico e o germânico, o que só é observável em alguns
casos de \emph{sandhi} já que o \emph{n} tende a se perder em posição final.
A semiconsoante *\pietrans{w} se transforma em uma fricativa labiodental /v/.

Para os grupos consonantais, sinaliza-se o resultado *\emph{dt}, *\emph{tt} >
\ocsex{st} (\cirilocsex{вєдѫ,vedõ} ``eu conduzo'', inf. \cirilocsex{вєсти,vesti}
< *\emph{ved-tēi}).
Por fim, a epêntese de uma consoante é bem atestada em dois contextos: a
palatalização de labiais produz um [l] epentético
(\emph{cf.}~\ref{sec:altmorfofon}) e em geral há uma dental epentética nos
grupos -\emph{sr}-, -\emph{zr}- (\cirilocsex{сєстра,sestra} ``irmã''
<*\pietrans{swesrA}, \cirilocsex{издрєшти,izdrešti} ``explicar'', de
\ocsex{iz-rešti}, \cirilocsex{въздрастъ,vъzdrastъ} ``idade'', de
\ocsex{vъz-rastъ}).
Contudo, os grupos -\emph{sr}-, -\emph{zr}- são viáveis
(\cirilocsex{срамъ,sramъ} ``humilhação'', \cirilocsex{zrakU} ``visão''): essa
aparente incoerência se explica pela cronologia relativa.
Os grupos -\emph{sr}- antigos produziram uma consoante epentética e a língua foi
dotada posteriormente novamente de grupos -\emph{sr}, -\emph{zr}- como
consequência de da metátese dos ditongos em líquida (\emph{cf.}~\ref{sec:silabasabertas}),
contexto em que não havia mais epêntese. A epêntese de [d] em um grupo -\emph{zr}-
secundário ainda é produtiva na língua, ainda que não seja automática: como
testemunho, há epêntese em palavra de empréstimo, como
\cirilocsex{издраиль,izdrail’ь} ``Israël'' e o alomorfe \ocsex{izd} da
preposição \ocsex{iz} diante de [r] em sandhi (\ocsex{iz-d-rǫky} ``da mão''
\passageme{L174}).


\subsection{Vocalismo}

O vocalismo foi consideravelmente mantido no eslavo comum, \emph{vide}~\autoref{tab:sistvoc}.

\begin{table}[!htb]
  \caption{Sistema vocálico do indo-europeu ao eslavo antigo}\label{tab:sistvoc}
  \begin{center}
    \begin{tabular}[c]{lllll}
      \toprule
      \emph{Vogais breves} & *ĭ > \emph{ь} & *ĕ > \emph{e} & *ă, ŏ > *ă >
      \emph{o} & *ŭ > \emph{ъ}\\
               & & \hspace{1em} > \emph{o}\slash{}\emph{-u̯} & & \\
      \midrule
      \emph{Vogais longas} & *ī > \emph{i} & *ē > \ocsex{E}
      [\ocsex{E}\textsubscript{1}] & *ā, *ō > *ā > \emph{a} & *ū > \emph{y}\\
      \midrule
      \emph{Ditongos} & &*ei̯ > [\emph{i}\textsubscript{1}] & *ai̯, *oi̯ > \ocsex{E}
      [\ocsex{E}\textsubscript{2}]/ \emph{i} [i\textsubscript{2}] & \\
                  &  & *eu̯ > \ocsex{ju} & *au̯, *ou̯ > \emph{u} & \\
      & \multicolumn{2}{l}{*ī̆n, *ī̆m, } & *an, *am, *on, *om > \ocsex{õ} & \\
      & \multicolumn{2}{l}{\hspace{0.5em}*en, *em > \ocsex{ẽ}} & *ŭn, *ŭm > \ocsex{õ} & \\
      \midrule
      \emph{Ressoantes}  & *r̥, *l̥ > \ocsex{rь lь}\slash\ocsex{rъ lъ} & & & \\
      \emph{vocálicas}  & *n̥, *m̥ > \ocsex{ьnV}, \ocsex{ьmV}\slash{}\ocsex{ęC} & & & \\
      \bottomrule
      \multicolumn{5}{l}{\small NB\@: após consoante palatalizada, *\pietrans{E} >
        \ocsex{E} > \emph{a}.}\\
    \end{tabular}
  \end{center}
\end{table}


As vogais *\emph{a} e *\emph{o} se fundiram, como em báltico e germânico, com
uma redistribuição de timbres em função da quantidade: i.e., *\emph{ă} e
*\emph{ŏ} > eslavo comum *\emph{ă} > eslavo \emph{o}, eslavo antigo \emph{o}
(inversamente em báltico e germânico, em que *\emph{ă} e *\emph{ŏ} > \emph{a}),
e *\emph{ā} e *\emph{ō} > eslavo comum *\emph{ā} > eslavo \emph{a} (germ.
*\emph{ā}, *\emph{ō} > \emph{o}, mas báltico oriental *\emph{ā} > \emph{o} e
*\emph{ō} > \glsxtrshort{lituano}~\emph{uo}); assim, \cirilocsex{дати,dati}
``dar'' < *\emph{dō-} < *\pietrans{deh3-} (\glsxtrshort{grego}~δίδωμι,
\glsxtrshort{lituano}~\emph{dúoti} ``dar'', \glsxtrshort{latim}~\emph{dōnum}
``dom''), \cirilocsex{домъ,domъ} ``casa'' < *\emph{domos}
(\glsxtrshort{grego}~δόμος, \glsxtrshort{latim}~\emph{domus}),
\cirilocsex{стати,stati} ``estar em pé'' < *\emph{stā-} < *\pietrans{steh2-}
(\glsxtrshort{grego}~ἵστᾱμι\slash{}ἵστημι, \glsxtrshort{latim}~\emph{stare},
\glsxtrshort{lituano}~\emph{stóti}), \cirilex{остръ} ``agudo'' <
*\pietrans{ăcro}- *\pietrans{h2ecro}- (\glsxtrshort{grego}~ἄκρος,
\glsxtrshort{latim}~\emph{Acer}).

A vogal *\emph{ū} se deslabializou em \emph{y}, que é uma vogal posterior fechada não
arredondada: \cirilocsex{сꙑнъ,synъ} ``filho'' < *\pietrans{suHnu-}
(\glsxtrshort{lituano}~\emph{sūnus}, \glsxtrshort{inglês}~\emph{son}). 

As vogais *\emph{ĕ} e *\emph{ē} resultam em, respectivamente, \emph{e} e
\ocsex{E} (quanto à realização fonética, \emph{cf.}~\ref{sec:invvoc}):
\cirilocsex{нєбо,nebo} ``céu'' < *\pietrans{nebhos} (\glsxtrshort{grego}~νέφος
``nuvem'', \glsxtrshort{latim}~\emph{nebula} ``idem''),
\cirilocsex{дѣꙗти,dEjati} ``agir'' < *\emph{dhē-} < *\pietrans{dheh1}
(\glsxtrshort{grego}~τίθημι, \glsxtrshort{lituano}~\emph{dëti} ``por''); mas
*\emph{ĕ} se labializa em o diante de *\emph{u̯}: \cirilocsex{новъ,novъ} ``novo''
< *\pietrans{newos} (\glsxtrshort{latim}~\emph{novus},
\glsxtrshort{grego}~νέ{(ϝ)}ος). 

As vogais *\emph{ĭ} e *\emph{ŭ} tornaram-se vogais reduzidas chamadas ``\emph{jers}''.
A tradição francesa as transcreve por \emph{ĭ} e \emph{ŭ} para o eslavo antigo,
embora eles não sejam realizados com os timbres [i] e [u]; essas vogais
reduzidas desapareceram nas línguas modernas e seu enfraquecimento já é visível
no eslavo antigo.
A transcrição por \grafema{ĭ} e \grafema{ŭ}, inadequada do ponto de vista
sincrônico, será abandonada em favor da notação com \grafema{ь} e \grafema{ъ}
que é a usada na escola anglo-saxã e alemã, sem transcrição.
Essa notação tem a vantagem de não falsear a interpretação e de melhor dar conta
das evoluções diacrônicas e das divergências dialetais no próprio eslavo
antigo.
Reservar-se-ão as grafias \emph{ĭ} e \emph{ŭ} para o estado do eslavo comum.
Assim \cirilocsex{ѥсмь,jesmь} ``eu sou'' <*\pietrans{h1es-mi}
(\glsxtrshort{sânscrito}~\emph{ásmi}, \glsxtrshort{grego}~εἰμί),
\cirilocsex{бъдѣти,bъdeti} ``velar'' <*\pietrans{bhudheh1}-
(\glsxtrshort{sânscrito}~\emph{budh-} ``ser desperto'',
\glsxtrshort{lituano}~\emph{budéti} ``velar'').

As ressoantes vocálicas se vocalizam com timbre [i] em geral, mas [u] em
contexto labial: *\pietrans{R} > *\emph{ĭr̥}, *\emph{ŭr} > \ocsex{ьr},
\ocsex{ъr}, *\emph{l̥} > *\emph{ĭl}, *\emph{ŭl} > \emph{ьl}, \ocsex{ъl},
*\emph{n̥} > *\emph{ĭn} > \ocsex{ьn}, *\emph{m̥} > *\emph{ĭm} > \ocsex{ъm},
conservados diante de vogal.
Quando a ressoante vocálica está diante de consoante, *\emph{r̥}e *\emph{l̥} se
tornam em eslavo antigo em \ocsex{rь}, \ocsex{lь} ou \ocsex{rъ}, \ocsex{lъ}
(\emph{cf.}~\ref{sec:silabasabertas} e~\ref{sec:invcons}) e *\emph{n̥}, *\emph{m̥}
pela vogal nasalizada \ocsex{ę}.

Os ditongos com segundo elemento nasal são ditongos apenas diante de consoante:
se eles estão diante de vogal, eles não formam ditongo e são conservados em
alteração. O mesmo ocorre com \ocsex{ьn} e \ocsex{ьm} que resultam de *\emph{n̥}
e *\emph{m̥}.
O resultado de *\emph{ĭn} e *\emph{ŭn} varia, com algumas palavras apresentando
a nasal esperada e outras um /i/, /y/ < *\emph{ū} que pressuporiam um resultado
*\emph{į}, *\emph{ų} em vogal nasal fechada, posteriormente desnasalizada.
Os detalhes desses desenvolvimentos continuam controversos, mas é claro que os
dois resultados são cronologicamente distintos: o resultado /i/, /y/ é
encontrado basicamente nos termos emprestados e não em palavras herdadas.


\subsection{Estrutura da sílaba}\label{sec:estruturasilabica}
\lipsum[2]
\subsubsection{Harmonia intra-silábica}
\lipsum[3]
\subsubsection{Lei das sílabas abertas}\label{sec:silabasabertas}
\lipsum[4]
\subsubsection{Resultados específicos em sílaba final}
\lipsum[5]

\subsection{Evolução posterior}
\lipsum[1]
